% --- Résumé ----------------------------------------------------------------------
\addchap{Résumé}

L'entreprise fictive \emph{InfraReady} souhaite standardiser la configuration de ses serveurs,
automatiser leur administration et les superviser. Le projet consistait à concevoir, déployer et
documenter un laboratoire représentatif : trois machines virtuelles sur un hyperviseur Proxmox VE,
dans un réseau isolé \code{192.168.50.0/24} — un contrôleur de domaine Windows Server 2025
(\textbf{WIN-DC01}, Active Directory et DNS), un serveur Ubuntu 24.04 hébergeant un service web
(\textbf{LINUX01}) et un serveur d'exploitation Ubuntu 24.04 portant Ansible, Prometheus et
Grafana (\textbf{OPS01}).

Le travail a été mené en huit phases (A à H) puis complété par deux séances de travaux pratiques.
Les installations sont automatiques, la configuration des serveurs est décrite par des playbooks
Ansible rejoués deux fois pour prouver leur idempotence, les secrets sont chiffrés par Ansible
Vault, et les métriques des trois serveurs sont collectées par Prometheus et présentées dans un
tableau de bord Grafana. Un dépôt GitHub versionne l'ensemble ; chaque modification passe par une
branche et une demande de fusion validée par un pipeline GitHub Actions (lint, syntaxe, règles
Prometheus, recherche de secrets).

La séance 6, consacrée à la maintenance opérationnelle et à la supervision, a ajouté un contrôle de
santé automatisé en lecture seule, une tâche planifiée qui teste le site web chaque minute, deux
alertes vérifiées par des pannes contrôlées, un tableau de bord d'exploitation, des quotas disque et
le diagnostic d'un incident réel : l'arrêt brutal du contrôleur de domaine à chaque arrêt de
l'hyperviseur.

\info{Résultats clés}{%
\begin{itemize}[leftmargin=1.2em]
  \item 3 serveurs opérationnels, domaine \code{infraready.test} fonctionnel, isolation réseau vérifiée ;
  \item 8 playbooks Ansible ; ceux de configuration sont idempotents (second passage \code{changed=0}) ;
  \item 4/4 cibles Prometheus collectées, 7 règles d'alerte, dashboard Grafana de 9 panneaux ;
  \item 16 demandes de fusion (PR), pipeline CI vert, gitleaks sans alerte sur tout l'historique ;
  \item plus de 100 preuves horodatées ; tous les échecs rencontrés sont documentés avec leur cause et leur correction.
\end{itemize}}

\attention{Points ouverts à la remise}{%
La correction de l'arrêt non propre du contrôleur de domaine (incident INC-2026-10-23-01) est
diagnostiquée, reproduite et préparée mais n'est pas appliquée : elle modifie le contrôleur de
domaine partagé et attend la validation de l'équipe. Les bilans individuels restent à rédiger par
chaque membre. L'emploi de Proxmox VE à la place de VMware Workstation doit être validé par le
formateur.}

\medskip
\noindent\textbf{Mots-clés :} Proxmox VE, Windows Server 2025, Active Directory, Ubuntu Server,
Ansible, idempotence, Ansible Vault, Prometheus, Grafana, node\_exporter, windows\_exporter,
GitHub Actions, maintenance, supervision, diagnostic d'incident.
